\chapter{Review of Related Literature}
\label{ch:literature}

\section{Treatment mechanisms and the meaning of model outputs}

Water quality is shaped by pollutant sources, transport, exposure, and the receiving environment. \citet{Madhav2019} organize pollutants by their sources and effects, while \citet{Lin2022} examine the different health consequences associated with polluted water. These works support a treatment assessment that considers several pollutant groups. They also explain why reducing a single effluent quantity cannot describe every environmental consequence of a treatment decision.

Biological treatment makes this assessment a problem of interacting transformations. \citet{Duan2022} review heterotrophic nitrification and its role within biological nitrogen removal. The review shows that a process label can contain several microbial pathways. \citet{Wentzel1992} connect nitrification, denitrification, and excess biological phosphorus removal through process conditions and model structure. For a surrogate, the implication is that nitrogen and phosphorus predictions belong to a coupled treatment response. Learning each reported total independently can omit relationships that matter when operating conditions change.

The ASM family provides a systematic framework for these relationships. \citet{Henze1999} describe the component and process structure of the model that includes biological phosphorus removal and denitrification by phosphorus-accumulating organisms. The collected models of \citet{Henze2006} place this development alongside other activated sludge formulations. \citet{Henze2008} provide the broader treatment principles needed to relate biological models to engineering design. These sources anchor the separation among the component state, the reaction model, and reported treatment indicators.

Carbon availability provides a concrete example of coupling. \citet{Guerrero2011} examine competition between phosphorus-accumulating organisms and denitrifiers under different carbon sources. Their work makes the identity and availability of carbon relevant to nutrient removal, rather than treating organic loading as a single interchangeable number. A component surrogate should therefore retain distinctions among fermentable material, readily available substrates, storage products, and biomass. The inclusion of interaction terms has a process rationale because the effect of one input can depend on another.

Oxygen supply has a similar role. \citet{StenstromSong1991} study the effect of oxygen transport limitations on nitrification. This distinguishes oxygen transfer from the biological capacity to consume oxygen. A model can include both and still give different responses when their relative limitations change. An educational explanation of aeration should therefore identify what is controlled, how oxygen enters the balance, and which reaction rates depend on dissolved oxygen.

The state of the biomass also depends on its operating history. \citet{Yilmaz2007} examine how alternating aerobic and anoxic or anaerobic conditions affect biomass activity during starvation. Their setting differs from a steady-state surrogate, but it identifies a limit of that surrogate. A mapping from current influent and operating conditions to one steady-state response does not represent all effects of prior exposure or prolonged substrate deprivation. Such evidence supports a precise boundary on claims about dynamic operation.

Model detail creates a calibration problem as well as a prediction problem. \citet{Machado2014} compare parameter identifiability with uncertainty in influent and operating conditions during full-scale model calibration. A close match to measurements does not uniquely identify every kinetic or stoichiometric parameter. This matters when interpreting a fitted surrogate. The surrogate approximates a response conditioned on the parameter values used to generate its training data. Its coefficients cannot be taken as newly identified biological constants.

\section{Development and calibration of activated sludge models}

The history of activated sludge modeling explains why a component response has an explicit scientific meaning. \citet{Marais1976} develop steady-state treatment relationships that connect biomass, substrate, and operating conditions. \citet{Dold1981} extend the representation toward a general process model. These works establish a distinction between an observed effluent total and the internal quantities used to calculate that total. A surrogate intended for component prediction inherits that distinction.

\citet{Henze1987} describe a general single-sludge formulation associated with Activated Sludge Model No.\ 1 (ASM1). Its process structure separates substrate use, biomass growth, decay, and nitrogen transformations. \citet{Gujer1999} develop Activated Sludge Model No.\ 3 (ASM3) with an explicit role for intracellular storage. The choice between these representations changes which intermediate states are available for prediction. A statistical model cannot infer a uniquely defined hidden storage state unless the supervised representation defines that state.

Nutrient-removal models extend the same principle. \citet{Barker1997} present a general biological nutrient-removal model, while \citet{Hu2007} develop a kinetic treatment of the linked processes. The phosphorus-removal formulation of \citet{Henze1999} adds another established component and process description. These sources support joint treatment of carbon, nitrogen, and phosphorus. They also show why a model name does not make every parameter or microbial pathway interchangeable.

\citet{Rieger2001} connect biological phosphorus removal to the storage-based formulation. \citet{Fall2015} modify an activated sludge model for low solids production. These contributions identify different purposes for model extension. Adding a process or changing a yield can alter the component relationships and the corresponding mass conservation conditions. An invariant operator must therefore be derived from the adopted stoichiometric matrix rather than borrowed from a model with similar names.

Model development has a practical history as well as a mathematical one. \citet{VanLoosdrecht2015} review the development and use of ASM1 over twenty-five years. \citet{Hauduc2013} examine the state of process knowledge, modeling concepts, and limitations. Their discussions support deliberate selection of a representation for a defined question. More states can describe additional mechanisms, but they also require additional information for calibration and assessment.

\citet{Petersen2003} review experimental designs for activated sludge calibration. \citet{Vanrolleghem1999} use respirometry to estimate model parameters and components. These studies connect a model quantity to the measurements that can support it. An oxygen-use observation can inform biological activity, but it is not a direct observation of every modeled component. The same issue arises when a surrogate is assessed against a few reported water quality measurements.

\citet{Meijer2004} examine the practical validation and calibration of a metabolic phosphorus-removal model. Their work distinguishes steady-state and dynamic information in testing a model. \citet{Nelson2009} analyze ASM1 mathematically. Together, these sources make clear that calibration, solution of the equations, and analysis of the resulting state are separate responsibilities. A good fit, a converged numerical solve, and an accepted physical state need not provide the same evidence.

The empirical biomass formula discussed by \citet{Hoover1952} supplies a useful historical example of constituent accounting. Biomass growth and decay redistribute organic matter and nitrogen, so a solids concentration cannot be treated as chemically empty material. \citet{TakacsVanrolleghem2006} examine elemental balances in activated sludge models and the role of omitted reaction products. Their contribution is especially relevant to projection. Mass conservation in an adopted lumped-component basis does not automatically establish complete elemental accounting for unrepresented water, carbon dioxide, or other boundary terms.

\section{Biological pathways and the boundaries of a model}

The functional role of a microorganism depends on both its energy source and its carbon source. Heterotrophic growth uses organic material, while nitrifying chemolithoautotrophs obtain energy from inorganic nitrogen transformations and use inorganic carbon for synthesis. \citet{Narayanan2019} review biological treatment and bioreactor design. \citet{Preena2021} review nitrification and denitrification in recirculating aquaculture. The latter setting differs from municipal activated sludge, but the distinction among electron donors, electron acceptors, and nitrogen forms remains useful.

Two-step nitrification separates ammonium oxidation from nitrite oxidation. It is an engineering representation of functional conversions rather than a universal rule about how many organisms perform them. \citet{VanKessel2015} establish complete nitrification within one microorganism. Their work explains why an explicit two-population model is a chosen process representation. It should not be interpreted as excluding every biological route that combines those functions.

\citet{Lu2014} review the microbial ecology of denitrification in biological wastewater treatment. Carbon availability, electron-acceptor conditions, and microbial populations jointly influence the nitrogen response. This supports retaining carbon and nitrogen components together when learning a surrogate. It also explains why the amount of reported nitrogen can change through conversion and separation even when the complete modeled nitrogen inventory satisfies mass conservation.

Anaerobic ammonium oxidation provides a different nitrogen pathway. \citet{Strous1998} study slowly growing anaerobic ammonium-oxidizing microorganisms. \citet{Dietl2015} investigate the hydrazine synthase complex, while \citet{Maalcke2016} examine the enzyme involved in nitrogen-gas production from hydrazine. These works clarify that a shared final nitrogen product does not make the underlying process identical to heterotrophic denitrification. The adopted reactor model does not contain an explicit anaerobic ammonium-oxidation population or its full kinetic pathway.

Temperature and treatment purpose create further boundaries. \citet{Vandekerckhove2018} investigate thermophilic nitrogen removal, and \citet{Xie2021} develop strategies for nitrogen conversion, recovery, and removal. \citet{WinklerStraka2019} review newer directions in biological nitrogen management. Their relevance is the need to identify which nitrogen transformation and engineering goal a model represents. A parameterized steady-state activated sludge surrogate does not acquire these additional capabilities through mass conservation and non-negativity alone.

\section{Suspended-growth and attached-growth treatment}

Suspended-growth and attached-growth systems organize biological activity differently. A completely mixed activated sludge reactor uses one bulk concentration vector for its control volume. A biofilm or packed treatment medium can contain spatial gradients and retained biomass at different locations. \citet{Wik2003} review trickling filters and biofilm modeling. \citet{Zhang1994} examine biofilm density, porosity, and pore structure. These contributions explain why transport and geometry can affect reaction rates even when the same nitrogen or organic-matter conversion is considered.

\citet{Machdar1997} develop a combined anaerobic and aerobic treatment setting for developing countries. \citet{Tawfik2006} examine combined anaerobic treatment and a down-flow sponge system. \citet{Tawfik2010} investigate the role of sponge volume in that downstream process. These studies connect treatment configuration and retained medium to the response being predicted. A medium volume or transport property is a different input from the hydraulic and aeration controls of a completely mixed reactor.

\citet{Mahmoud2010} study a naturally ventilated bio-tower for organic and nutrient treatment, and \citet{Mahmoud2011} examine a sponge reactor for municipal wastewater post-treatment. \citet{Hatamoto2018} relate sponge-reactor structure to process function and microbial communities. These contributions establish attached-growth treatment as a relevant engineering context. They do not supply additional evaluated cases for the present activated sludge benchmark.

\citet{Hellal2021} simulate a passively aerated biological filter, while \citet{Liang2021} model a three-stage biological trickling filter. \citet{Laine1999} examine dynamic nitrification in a low-loaded trickling filter. Their different configurations show why stage structure, spatial transport, and time dependence need to be identified before model responses are compared. A surrogate can be fast while approximating a different physical problem from the one intended.

Nutrient form also matters in these systems. \citet{Kasi2011} model dissolved organic nitrogen in a two-stage trickling filter, and \citet{Simsek2012} examine its fate. \citet{BressaniRibeiro2021} investigate inorganic carbon limitation during nitrogen conversions in sponge-bed treatment. These works connect prediction targets to constituent forms and growth requirements. They reinforce the need to retain the component basis rather than treating one measured nitrogen total as a complete process state.

\citet{DiezMontero2019} assess a trickling-filter upgrade for nutrient removal. \citet{Luan2023} examine nitrification and denitrification in a rotating self-aerated biofilm reactor. Both illustrate the engineering importance of configuration-specific process relationships. A statistical predictor trained for one hydraulic and biological arrangement requires new evidence before transfer to another.

The resource setting is also relevant. \citet{Tyagi2021} review the energy-saving motivation for sponge treatment, and \citet{Nasr2022} discuss its decentralized applications. \citet{Rapi2021} apply biological sponge treatment to palm oil mill effluent. These sources broaden the context of treatment selection while keeping the wastewater type explicit. Industrial and municipal influents can differ in their constituent distributions and treatment demands.

\citet{Mahmoud2018} consider post-treatment of anaerobic effluent containing a specific organic contaminant and heavy metals. Their setting illustrates a further limitation of a declared component model. The presence of metal-hydroxide and metal-phosphate coordinates does not make an activated sludge surrogate a general predictor of toxic-metal removal or every organic contaminant. The represented species and reactions determine that capability.

Transport modeling can join reaction kinetics with flow and diffusion. \citet{SadinoRiquelme2023} review how biological kinetics are integrated with computational fluid dynamics. The concern is the spatial location and transport of reacting material. The current standalone reactor and connected plant use their declared completely mixed stages and clarifier representation. Spatial treatment models provide a useful boundary comparison and a possible extension, rather than a replacement for those reference equations.

\section{Why treatment analysis calls for surrogate models}

Treatment design often contains competing objectives and coupled decisions. \citet{PadronPaez2020} use multiobjective optimization to address sustainable wastewater treatment plant design. \citet{Aboagye2021} link wastewater network design to optimization and sustainability assessment. Their contributions show why treatment evaluation involves repeated comparisons among alternatives. The computational burden belongs to the decision problem as well as to an individual simulation.

Sensitivity analysis gives another reason for repeated evaluation. \citet{Araujo2013} examine optimal operation and the selection of economic controlled variables in an activated sludge model. The work connects operating sensitivities to the quantities that are useful for control. A fast approximation can assist this analysis, but it must retain the correct input definitions and reporting basis. A sensitivity with respect to reactor volume is different from a sensitivity with respect to a residence time that uses a recycle-adjusted flow.

The review of \citet{BhosekarIerapetritou2018} treats surrogate construction, feasibility analysis, and optimization as connected tasks. It emphasizes the importance of sampling, model selection, and validation for the intended application. \citet{McBride2019} place surrogate modeling within chemical process engineering and distinguish different ways of approximating an expensive response. These reviews support a broad comparison of learning methods. They do not imply that one statistical family is best for every response or data regime.

\citet{Fang2011} develop a simulation-based integrated approach to biological nutrient-removal optimization in a full-scale wastewater treatment plant. Their contribution links operating choices to the coupled nitrogen and phosphorus response through a process model. It also identifies the reference calculation that a statistical surrogate seeks to reduce. The surrogate can approximate repeated evaluations, but the original component and treatment relationships remain the basis for assessing an operating choice.

Wastewater resource recovery provides a direct application. \citet{Durkin2024} combine process surrogates with optimization of resource recovery from wastewater. The statistical approximation is useful because it participates in a larger constrained decision problem. \citet{He2023} use a Kriging surrogate in a papermaking wastewater treatment application with environmental objectives. Both studies connect faster response evaluation to engineering choice. Their relevance is methodological even when the treatment processes and outputs differ from an activated sludge component surrogate.

Recent network work extends this concern across treatment stages. \citet{Tang2026} develop surrogate-based design for a multistage coal chemical wastewater treatment network. A network surrogate must preserve the meaning of the connections among units as well as the response of each unit. \citet{Tsochatzidi2025} develop a surrogate optimization framework for pharmaceutical process systems. That work is outside wastewater treatment, but it illustrates the same distinction between approximation of process responses and optimization over the resulting model. Transfer to activated sludge requires the treatment-specific component and hydraulic constraints.

Reduced models can also be chosen around quantities that matter to operation. \citet{StrausSkogestad2018} generate surrogates using self-optimizing variables. Their contribution concerns the connection between reduced representation and decision purpose. A surrogate that is adequate for one operating objective can be inadequate for another. This supports retaining complete component predictions when the later objective may combine nutrients, organic matter, solids, and resource use.

\section{From local approximations to nonlinear data-driven prediction}

Local approximation has a clear engineering role. \citet{Li2020} use fuzzy clustering in dissolved oxygen model predictive control, while \citet{Yang2014} develop fuzzy model-based predictive control for the activated sludge process. Such approaches organize nonlinear behavior through local models or operating regimes. Their usefulness depends on how well the local representation covers the conditions encountered in operation. They provide an alternative to fitting one highly flexible global predictor.

Polynomial and response surface methods make nonlinear terms explicit. \citet{Kim2009} develop a dual optimization strategy for nitrogen and phosphorus removal. \citet{Lim2011} examine operating conditions for organic matter and nitrogen removal in a membrane bioreactor. \citet{Ahn2014} compare modeled and experimental operating factors for livestock wastewater treatment. These works demonstrate why transparent approximations remain useful in treatment analysis. The terms have an inspectable structure, but the selected polynomial degree and sampled domain limit the behaviors they can represent.

\citet{Nazlabadi2021} review the application of response surface methodology to biological wastewater treatment. The review places polynomial approximations within a wider design and analysis practice. A low-order surface can reveal curvature and input interaction without reproducing every mechanistic detail. It should therefore be assessed as an approximation over a stated domain. This perspective supports second-order regression while preventing a claim that a quadratic structure is a complete reaction model.

Partial least squares provides another transparent starting point. \citet{Wold2001} explain how latent variables relate predictors to responses in chemometric regression. \citet{Khatri2020} review its use in water quality analysis. The method is relevant when influent variables are correlated and when several outputs are modeled together. Its basic linear relationship still depends on whether the predictor representation captures the nonlinear behavior of interest.

Extensions address this limitation in different ways. \citet{Woo2009} use kernel partial least squares to estimate process variables in industrial wastewater treatment. \citet{Liu2021} use a dynamic concurrent kernel approach for wastewater process monitoring. \citet{Yang2021} embed a relevance vector model within a partial least squares framework for adaptive effluent prediction. These contributions show that latent structure, nonlinearity, and changing process behavior can be combined. They also separate dynamic monitoring from the steady-state input--output task considered here.

The wider evidence on wastewater data use is reviewed by \citet{Newhart2019}. Their discussion makes data quality, process knowledge, and validation important parts of performance analysis. \citet{BaAlawi2021} address influent sensor validation through stacked denoising autoencoders. Their contribution concerns the credibility of the inputs before a prediction is made. An accurate model cannot compensate for every error in a sensor signal or a mislabeled input quantity.

\citet{Bagherzadeh2021} examine total nitrogen prediction and feature selection in a wastewater treatment plant. \citet{SadriMoghaddamMesghali2023} develop a hybrid ensemble for total nitrogen prediction in a combined trickling filter and activated sludge system. These studies demonstrate the value of statistical prediction for an operational indicator. Their output target also clarifies a limitation. Evidence for one total nitrogen model does not establish mass conservation or non-negativity of the complete component vector.

Recent studies connect prediction to other practical indicators. \citet{Wang2025} compare interpretable learning models with a mechanistic model for effluent nitrogen prediction. \citet{Yang2025} examine waste activated sludge generation for operational efficiency. \citet{Xiao2026} combine mechanistic dynamics and an interpretable surrogate in activated sludge management. Together, these applications broaden the role of learning from effluent estimation to sludge and operating analysis. Their different targets require careful comparison because an error measure over one indicator is not an error measure over the full state.

\citet{Asadi2017} use data mining in aeration optimization, and \citet{Mao2024} use artificial intelligence to balance effluent quality and aeration energy. These works give a concrete decision context for prediction. Oxygen delivery can support treatment while consuming resources. A useful operating surrogate must represent this relationship well enough at the selected decision, rather than only on average over randomly held-out inputs.

\section{Prediction targets and information available to a decision}

Wastewater prediction studies often address a reported quantity rather than a complete mechanistic state. \citet{Ching2022} develop a soft sensor for five-day biochemical oxygen demand. \citet{ManavDemir2024} compare learning methods for nutrient-removal effluent parameters. These applications demonstrate the relevance of prediction to treatment monitoring. Their supervised target determines what an error score can establish. Agreement with one reported quantity does not establish mass conservation or non-negativity across the full component vector.

\citet{Ekinci2023} examine feature selection for sludge-production prediction. \citet{Wang2024} predict COD component parameters for activated sludge modeling. These contributions concern different output roles. A sludge production quantity, a wastewater fractionation parameter, and an effluent concentration are not interchangeable labels. The predictor inputs must also match the information available when each quantity is needed.

\citet{Poch1993} provide an observational treatment-plant dataset for classifying plant state and detecting faults. It contains measurements from different stages and performance quantities. A measurement can be a legitimate input for an operational diagnosis while remaining unavailable for prospective design or steady-state prediction at an untested setting. The dataset consequently helps explain information availability. It is not the training collection used for the current reactor or plant surrogates.

\citet{Caro2024} describe automated acquisition of simulator responses. Their contribution is relevant to the organization of paired model-generated observations. Such acquisition does not establish agreement with field measurements or eliminate the need for independent physical and numerical checks. The review of \citet{Sundui2021} places learning applications across biological wastewater tasks and similarly supports identifying the data source and intended output before comparing models.

Full-scale mechanistic applications provide complementary evidence. \citet{Wu2016} simulate and optimize a coking wastewater process with activated sludge models. \citet{SadriMoghaddam2021} address calibration of a full-scale plant model, and \citet{Jasim2020} develop a treatment-plant design model. These studies connect a model to a particular process setting. Their calibration does not remove the need to assess a surrogate separately on its sampled operating domain.

\citet{Szelag2022} examine nutrient removal and energy consumption under uncertainty. \citet{Nazif2023} formulate practical operating analysis with attention to influent characteristics. These contributions strengthen the decision context of model assessment. The operating response depends on the influent and parameter assumptions, so a selected decision should be evaluated at the stated conditions. Physical constraints do not remove every uncertainty in those conditions.

\section{Learning families and what each comparison can reveal}

Different learning methods offer different ways to approximate the same response. \citet{Hastie2009} provide the statistical foundations for regularization, ensembles, kernels, and model assessment. \citet{ShwartzZiv2022} examine tabular data and caution against assuming that deep learning is necessary for every structured dataset. \citet{Borisov2024} review deep neural networks for tabular learning. These contributions motivate a comparison that includes simple, nonlinear, and neural methods under the same prediction target.

Tree ensembles form one important group. \citet{Breiman2001} develop random forests by aggregating randomized trees. Averaging can reduce the instability of a single tree while retaining nonlinear input partitioning. \citet{Geurts2006} introduce extremely randomized trees with stronger randomization of split choices. Their different construction makes both methods useful in a benchmark. Neither construction imposes a coupled material balance across outputs unless that relation is added explicitly.

Boosting builds an ensemble sequentially. \citet{Drucker1997} adapt boosting to regression by increasing attention to difficult observations. \citet{ChenGuestrin2016} develop extreme gradient boosting (XGBoost) with regularized tree boosting and efficient treatment of large datasets. \citet{Ke2017} develop the light gradient boosting machine (LightGBM) with methods for efficient gradient boosting. \citet{Prokhorenkova2018} develop CatBoost with an ordered approach that addresses prediction shift and categorical variables. The component data considered here are numeric, but these algorithms still provide different nonlinear fitting procedures. Their names do not establish physical admissibility or superiority in advance.

Neighbor and kernel methods give complementary comparisons. \citet{Mack1981} study the local properties of nearest-neighbor regression estimates. The prediction is strongly connected to the distribution of nearby training inputs. This makes domain distance relevant when the requested input moves outside the sampled range. \citet{SmolaScholkopf2004} explain support vector regression and its control of model complexity through an error-insensitive fitting objective and regularization. A kernel can represent nonlinear relationships, but the learned relationship still requires separate assessment in an extrapolation regime.

Multiresponse regularization addresses structure across outputs. \citet{Zou2005} combine shrinkage and variable selection through the elastic net. \citet{Yuan2006} develop regression selection for grouped variables. \citet{Argyriou2008} treat feature learning across several tasks. \citet{Liu2009} provide blockwise procedures for the multi-task lasso, and \citet{Simon2013} develop blockwise descent for group-penalized multiresponse regression. These works justify comparisons that can share information across component outputs. A shared sparsity pattern is a statistical relation and should not be mistaken for stoichiometric mass conservation.

Neural regression offers a flexible function representation. \citet{Rumelhart1986} develop learning through back propagation, and \citet{Hornik1989} establish approximation properties of multilayer feedforward networks. Such approximation results concern representational capacity under stated conditions. They do not provide an automatic guarantee about finite-data accuracy, optimization success, extrapolation, or material balance. \citet{Curteanu2011} discuss network topology in chemical applications, making architecture selection part of the modeling problem.

\citet{ArikPfister2021} introduce TabNet with sequential attention over tabular inputs. Its feature masks can support inspection of which inputs are used. Attention-based selection still differs from a named algebraic coefficient or a component balance for mass conservation. It is useful to compare such a model with an explicit second-order regression because the two forms of interpretability answer different questions.

\section{Learning behavior, sampling, and fair assessment}

Training error alone gives incomplete evidence of approximation quality. \citet{Geman1992} analyze the bias and variance dilemma in neural learning. A flexible model can reduce systematic approximation error while becoming more sensitive to the training sample. \citet{Belkin2019} show that modern learning behavior can be more complex than a single classical tradeoff curve. These contributions support direct learning-curve assessment rather than assigning a fixed sample-efficiency ranking to a model family.

\citet{VieringLoog2023} review the shapes of learning curves and the conditions that influence them. The curve relates available training data to predictive performance under an explicit evaluation design. Its interpretation depends on the error scale, spacing of sample sizes, and variability of the data partitions. A normalized area under such a curve is a useful summary only when those choices are common across the compared models.

Training inputs must also cover the intended process domain. \citet{McKay1979} compare sampling procedures and motivate stratified coverage of input ranges. Latin hypercube sampling can spread observations across each selected input range without enumerating every combination. It does not guarantee equal coverage of every joint response regime, nor does it make independently sampled influent components representative of every real wastewater composition. These limitations need to remain visible when simulated data define the benchmark.

Active learning raises a related question about the cost of obtaining informative targets. \citet{Guyon2011} organize an active learning challenge in which labels are limited. \citet{Cawley2011} provide simple regularized baseline strategies and emphasize the value of comparing an elaborate selection method with simpler alternatives. Their contribution is relevant to simulator data because each reference response also has a cost. The present learning-curve design measures dependence on available labels. It does not claim to implement or validate an active selection strategy.

Hyperparameter selection is part of model fitting. \citet{Bergstra2011} develop model-based methods for searching hyperparameter spaces. \citet{Akiba2019} provide a framework for organized hyperparameter optimization. These methods can improve the use of a search budget, but they require separation between the observations used to select settings and the observations used to assess the selected model. The same requirement applies when choosing a projection threshold or a trust criterion.

\citet{Kaufman2012} formalize leakage in data mining and explain how information that should be unavailable can enter the fitted model. In this setting, scaling all observations before cross-validation would allow validation inputs to influence training transformations. Selecting hyperparameters against an outer assessment set would allow the same set to guide and judge the model. Nested partitioning and training-only transformations are therefore part of the evidence design rather than minor computational details.

\citet{Pinheiro2025} examine how feature scaling affects learning tasks. Their contribution supports an explicit account of which methods use scaling and how those scales are fitted. Scaling also changes the numerical meaning of a coefficient. A standardized second-order term must be transformed back to physical input units before an engineering interpretation is attempted.

Comparison must account for variation across assessment problems. \citet{Demsar2006} develop statistical comparisons for classifiers over multiple datasets. The present application concerns regression and repeated data partitions, so its recommendations require adaptation to the unit of independent evidence. \citet{Shevchenko2024} examine benchmarking variability in recommender systems. Their setting is also different, but the concern transfers. Reported rankings can depend on splits and evaluation choices. Paired model comparisons and transparent variability are more defensible than a universal ranking inferred from one split.

Error normalization deserves its own explanation. \citet{Tofallis2015} examine relative prediction accuracy and the problems created by some relative measures. Near-zero component concentrations make ordinary percentage error unstable. A small absolute error can become an enormous percentage or remain undefined at zero. Componentwise normalization using a declared training reference provides a different comparison basis. It does not eliminate the need to report physically meaningful absolute errors and constraint violations.

\citet{Shahbazi2026} address imbalanced regression through data-level and algorithm-level balancing. This concern is relevant when rare operating conditions or low residual concentrations matter more to a decision than their frequency suggests. An average error over the sampled domain can conceal such behavior. Component-resolved errors and selected-decision evaluation are therefore needed alongside aggregate scores.

\section{Physical constraints in scientific machine learning}

Scientific machine learning combines data-based approximation with knowledge of physical structure. \citet{Karniadakis2021} review this combination across several forms of physics-informed learning. \citet{Kashinath2021} discuss applications in weather and climate modeling. These reviews show that physics can enter through inputs, architecture, training, or outputs. The practical question is which relation is enforced, at what stage, and under what conditions.

\citet{Raissi2019} develop physics-informed neural networks that include differential equation residuals in a learning objective. This approach can use governing equations even when direct response data are limited. A finite residual penalty still permits a tradeoff with the data-fitting term and other penalties. Exact mass conservation of every returned component state requires an additional structural argument or an explicit feasibility procedure.

Wastewater-specific work makes the same distinction useful. \citet{Li2024} apply physics-informed learning to unit operations and processes, including an activated sludge reactor. The governing-equation information supports the approximation task. Its inclusion in training does not automatically establish every plant-wide hydraulic or clarifier balance. Those relations must be defined at the level where the coupled outputs are returned.

Architectural enforcement supplies one stronger route. \citet{Beucler2021} impose analytic constraints in neural models by arranging the output relationship so that the required equalities hold. Some outputs can be reconstructed from others to complete a constrained state. This depends on which quantities are independent and whether the reconstructed quantities remain within their allowed bounds. Equality completion alone does not settle non-negativity.

\citet{SturmWexler2020} formulate a approach for mass conservation for photochemical computations. \citet{SturmWexler2022} embed atom balance in a neural architecture through the stoichiometric relationship. These contributions connect learned transformations to reaction structure. \citet{Sturm2023} use dimensionality reduction with mass conservation for transport of organic aerosol. Their scaling can also maintain non-negative concentrations in the reduced representation. This provides a precedent for addressing mass conservation and non-negativity together while reducing the number of transported quantities. Its transport setting differs from steady-state effluent projection. In activated sludge, an analogous invariant construction must respect the native component basis and the external forcing terms.

Reconciliation provides another route from a statistical estimate to a consistent state. \citet{Hansen2024} develop learning models that respect mass conservation through a probabilistic reconciliation perspective. Mass conservation supplies information that can be combined with the learned prediction. This differs from assuming that a low training loss has already made the prediction consistent with mass conservation. The connection to activated sludge is the availability of known linear inventory relations. The probabilistic interpretation remains distinct from a deterministic projection that returns one state without a calibrated uncertainty distribution.

\citet{Ruckstuhl2021} learn a projection for a data assimilation procedure from the difference between an unconstrained analysis and a more expensive constrained analysis. The training target contains information about the projection that the physical constraints require. This is a different task from learning a reactor effluent directly from influent and operating conditions. It also illustrates a separate computational choice. One can solve a constrained problem at prediction time or train another model to approximate its projection. If the projection is approximated, the returned state still requires a feasibility check. The information carried by constrained training examples should not be confused with exact enforcement of the constraints on every new prediction.

\citet{Geiss2022} address mass conservation when a neural model converts coarse atmospheric chemistry fields to a finer spatial grid. Their formulation accounts for unequal grid-cell areas and consistency between coarse inputs and fine outputs. It also treats the normalization of quantities whose distributions span large concentration ranges. This application clarifies why a physical operator must match the representation on which it acts. A spatial average needs area weights, while a treatment inventory needs native-unit constituent weights and a specified boundary. Directly borrowing a layer for mass conservation without transforming those weights would enforce a different relation. The wastewater application therefore adopts the principle of explicit physical accounting while deriving its own component operator.

\citet{Harder2022} emulate aerosol microphysics by learning changes in state variables from mechanistic input--output pairs. They compare physical information in soft loss terms with completion and projection layers, and examine non-negativity enforcement. Their tendency target gives an additional lesson. A negative change can be physically meaningful when it represents consumption, even though the resulting concentration must remain non-negative. The sign requirement belongs to the appropriate quantity. In activated sludge, driver coefficients and reaction-induced changes can have either sign, while the returned component state is subject to non-negativity. These distinctions help specify what a constrained surrogate actually guarantees across different scientific tasks.

Output non-negativity is another physical requirement. \citet{Haddad2010} provide a mathematical treatment of non-negative and compartmental dynamical systems. In a concentration model, non-negativity belongs to the state description itself. Clipping a negative value can satisfy this bound while changing the total represented material. A projection that enforces mass conservation alone can redistribute material into a negative component. Simultaneous enforcement must therefore handle both requirements together.

Hard constraint layers also raise questions about expressive capacity and uncertainty. \citet{Min2024} study an input-dependent enforcement layer for affine and convex output constraints with neural approximation guarantees. These guarantees depend on the stated architecture and assumptions. They do not establish the approximation capacity of a finite second-order regression. \citet{Utkarsh2025} develop a probabilistic projection that uses prediction-variance information to impose hard constraints and assess uncertainty. These methods incorporate feasibility into the learning architecture or prediction distribution. The deterministic projections used here have different weighting choices and do not inherit a probabilistic uncertainty guarantee.

\citet{Lakshminarayanan2017} develop deep ensembles for predictive uncertainty estimation. Spread among models can help describe uncertainty in a prediction. It does not certify a component balance or establish the validity of an extrapolated state. Likewise, a deterministic projection does not by itself produce a calibrated uncertainty distribution. Physical admissibility and uncertainty assessment are separate parts of model credibility.

\section{Projection and the geometry of admissible states}

Projection restricts a completed prediction to a set defined by physical relations. \citet{SturmSilva2024} use a species-weighted projection for mass conservation through atom balances in atmospheric composition. The weighting matters because species can differ greatly in concentration. An equal absolute adjustment can be small for a major component and large for a trace component. This motivates an explicit choice of projection geometry rather than treating every projection for mass conservation as equivalent.

\citet{Valente2025} project predictions onto physical manifolds defined by mass conservation laws. Their work gives output projection a general geometric interpretation. A predicted point is moved to an admissible set according to a selected notion of distance. The interpretation of the projection depends on the distance, the set, and the existence of a feasible state. A closer point in one weighted distance need not have a smaller error in another evaluation metric.

\citet{KircherVotsmeier2025} combine mass conservation and non-negativity for mechanistic kinetics surrogates. The procedure uses a positive provisional prediction to define relative weighting, obtains a candidate satisfying mass conservation, and enforces non-negativity through interpolation with an admissible reference when needed. This construction can be applied after different statistical predictors have been fitted. It is valuable for isolating the effect of physical enforcement from the effect of changing the learner.

The reference state has a practical meaning. If the influent is non-negative and has the required inventories used for mass conservation, it lies in the feasible set for the standalone reactor invariant constraints. Interpolating toward this state can enforce non-negativity without changing those inventories. The resulting state is physically admissible under the imposed constraints. It need not reproduce the reference reactor kinetics. This distinction prevents a guarantee about mass conservation from becoming an unsupported claim about treatment accuracy.

The geometry also limits claims about accuracy improvement. Projection onto a closed convex set in a fixed Euclidean norm has useful distance properties when the true state lies in that set \citep{Boyd2004}. A relative-weighted candidate satisfying mass conservation followed by boundary interpolation is a different map. Its weights can depend on the prediction, and its final interpolation is not generally the nearest feasible point under the reported error metric. Paired empirical assessment is therefore necessary even when feasibility follows from the construction.

\section{Interpretable constrained regression}

\citet{Rudin2019} argue for interpretable models in high-consequence settings when such models are available. The argument concerns the ability to inspect the model itself, rather than relying only on an explanation added afterward. For treatment surrogates, an explicit regression can identify operating terms, loading terms, nonlinear terms, and component coupling. This creates a structure that an engineer can inspect in relation to the process representation.

\citet{Brunton2016} identify governing relationships through sparse candidate functions in dynamical systems. Their contribution shows how named terms can connect learned structure to scientific interpretation. A second-order treatment surrogate uses named terms for a different purpose. It approximates a steady-state response and does not claim to discover the biological reaction laws. Similar predictions can arise from different coefficient combinations when inputs are correlated or when several blocks compensate for one another.

\citet{Shen2023} discuss differentiable modeling that joins physical and learned components. This supports inspection of the entire prediction chain. For a coupled regression, the coefficients of the driver are transformed by the coupling solution. For a projected model, the final output also depends on the active constraints. An explanation that ignores these transformations describes only part of the returned prediction.

The published formulation of \citet{Selerio2026} supplies a specific activated sludge example. It combines an invariant-aware second-order regression with a checked output procedure that imposes mass conservation and component non-negativity. The important theoretical distinction is between the unconstrained regression response, the fitted state variables used during training, and the final deployed state. Soft invariant-aware fitting guides the regression. The checked projection establishes the stated final-state constraints. These are different mechanisms and require different diagnostics.

Numerical optimization underlies both mechanisms. \citet{Boyd2004} provide the convex optimization foundations for affine constraints and quadratic objectives. \citet{Razaviyayn2013} analyze block successive minimization under specified regularity conditions. A regression objective that jointly fits coupling and state variables can be nonconvex even when individual updates are convex. Decreasing the objective at each accepted update does not establish a globally optimal fit. Conditioning restrictions also need their own treatment because they are not ordinary box constraints.

\citet{Stellato2018} present an operator-splitting approach for quadratic programs, and \citet{Stellato2020} give its fuller mathematical and computational development. \citet{HuangfuHall2018} develop a parallel dual revised simplex method for linear programming. These works support the numerical treatment of regression subproblems and output projections. Their role is to solve specified optimization problems. Numerical termination still needs explicit residual and bound checks before a predicted state is accepted.

\section{Connected plants and the move from prediction to decisions}

A connected activated sludge plant contains more physical structure than one reactor. \citet{HartelPopel1992} develop a dynamic secondary clarifier model with sludge thickening. \citet{Takacs1991} model clarification and thickening through a settling relationship. \citet{JeppssonDiehl1996} examine the propagation of biological components through the secondary clarifier. These contributions explain why clarification must preserve component relationships and why overflow, underflow, and stored solids cannot be predicted as unrelated quantities.

\citet{Alex2008} define a benchmark simulation setting that connects biological treatment, clarification, and recirculation. Its value is a common process context for comparing operating strategies. A plant with an extended component basis remains distinct from that benchmark, but can use the same discipline of declared hydraulics, unit connections, and assessment conditions. In particular, a recycle stream affects both concentration mixing and the flow passing through each reactor.

Solids retention adds another connection. \citet{Leu2012} examine the implications of long solids retention time for energy use, effluent quality, and stability. The retained solids depend on both reactor and clarifier inventories and on external solids losses. A surrogate that predicts effluent quality accurately but neglects inventory can therefore misrepresent a solids-age constraint. \citet{Miederer2025} develop an energy management model for wastewater treatment plants, further connecting operating choices to resource use.

The \citet{USEPA2009} review provides engineering guidance for nutrient control technologies. Its role is to identify treatment considerations and design relationships. Site-specific effluent limits and safety thresholds still belong to the declared optimization problem. They should be supplied as explicit inputs rather than presented as universal regulatory limits.

Control studies show how a prediction model enters operation. \citet{Sadeghassadi2018} use neural models for setpoint design and predictive control. \citet{Bernardelli2020} combine a learning model with real-time predictive control of a wastewater treatment plant. \citet{Ren2026} study ammonium-based aeration scheduling under influent disturbances. \citet{Croll2023} review reinforcement learning for wastewater process control. These contributions demonstrate a range of operating applications. Their dynamic decisions remain distinct from a steady-state operating search, which evaluates equilibrium responses under fixed controls.

Joint output projection offers a way to preserve the specified connections in a statistical plant model. Mixer equations, reactor invariants, clarifier outlet balances, and inventory bounds can be imposed together. At fixed controls, the relevant relations can form a convex quadratic projection problem \citep{Boyd2004,Stellato2020}. The projection can remove contradictions among separately predicted unit responses. It does not impose every kinetic rate equation or every layerwise settling equation of the mechanistic plant.

This creates an optimization problem with a nested prediction procedure. \citet{Colson2007} review bilevel optimization, which helps explain the distinction between an outer control choice and an inner constrained response calculation. The plant projection has a strongly convex inner objective when the weights are positive and the feasible set is nonempty. \citet{Tondel2003} analyze parametric quadratic programming and its relation to piecewise solution structure. Their results motivate attention to active constraints, but an operating problem with control-dependent balance coefficients need not have the same globally piecewise-affine structure.

\section{Water networks and the distinction between topology and operation}

Water-network design separates the origins, combinations, and destinations of streams. A source supplies water with specified quantities and quality. A mixer combines compatible streams. A sink accepts water for a use or discharge requirement. \citet{Galan1998} optimize distributed wastewater treatment networks, and \citet{Galan1999} examine design and synthesis strategies for those networks. Their contribution concerns the choice of connected treatment and routing alternatives. The present connected plant uses a fixed topology and varies declared operating controls.

Regeneration and constituent-specific treatment complicate those choices. \citet{GuelliSouza2011} consider differentiated regeneration of contaminants in chemical-industry water reuse. \citet{YangNetwork2014} include wastewater regeneration models in network optimization. The treatment response influences which stream can reach which destination. These works show why a constant removal fraction can be an insufficient description when concentrations and operating conditions change.

Industry-specific applications make the boundary explicit. \citet{Gutterres2010} examine water reuse in tannery operations, \citet{Hansen2018} minimize water and wastewater in a petrochemical setting, and \citet{Xu2015} integrate a yeast-production water system. The application determines the relevant contaminants, unit operations, and water-use requirements. Their network methods are useful conceptual precedents, while their process states are not the activated sludge component targets.

\citet{RubioCastro2013} develop global optimization for property-based interplant water integration. Their result concerns the specified mathematical formulation and the conditions required by its optimization method. \citet{McCormick1976} provide convex underestimating relations for nonconvex factorable expressions. Such relations can support bounds on an optimization problem. They do not establish that a projected surrogate response satisfies every original nonlinear process equation.

\citet{Tosarkani2020} design a wastewater treatment network under uncertainty through a multiobjective robust formulation. \citet{Ye2019} use a hybrid particle-swarm approach for network planning. These contributions address uncertainty and search over a planning problem. The numerical route and stopping evidence must be distinguished from a global optimality certificate. A faster operating search in a fixed plant is a different contribution from selecting a globally optimal treatment network.

\citet{DamalerioRemoval2022} connect phosphorus-removal priorities with treatment-configuration and cost objectives. \citet{Selerio2022} develop regression-based optimization for urban water eutrophication abatement. Both connect response approximation to a declared decision purpose. Their relevance is the organized treatment of several engineering quantities. The current study retains a component-based activated sludge response, projection for mass conservation and non-negativity, and independent verification of selected controls.

\section{Approximation and the structure of an optimization claim}

\citet{KusiakWei2013} examine data-based optimization of the activated sludge process. Their contribution connects a learned response to operating selection. \citet{Niu2022} combine deep learning and multiobjective dynamic wastewater optimization. These studies illustrate different response representations and search settings. A steady-state operating comparison should identify its time basis and objective rather than treating all learning-based optimization as one task.

\citet{Jana2022} combine neural and support-vector regression methods in detergent-industry effluent optimization. \citet{Waqas2022} model membrane permeability for a rotating biological contactor, while \citet{Li2022} join response-surface and neural methods for membrane fabrication. These applications show how statistical approximation can enter an engineering search. They also show why evidence for one target cannot be transferred directly to a complete reactor component state.

\citet{Shojaei2021} discuss the optimization of conditions in wastewater degradation. \citet{Poh2016} review mechanistic and metaheuristic optimization in anaerobic digestion. The process models and decision variables differ from those of activated sludge, but the common issue is the relation between an inexpensive response calculation and a constrained operating choice. Model accuracy, search completeness, and the chosen priorities remain distinct.

\citet{JainKar2017} examine nonconvex optimization in machine learning. A model can be linear in its fitted coefficients while remaining nonlinear in its operating inputs. A convex projection at fixed inputs can also be embedded in a nonconvex outer search. This distinction is essential to the present contribution. Neither transparent regression terms nor an exactly solved projection alone establish a globally optimal operating decision.

\section{Managing approximation during operating searches}

An optimizer can exploit approximation error in a favorable region. \citet{Alexandrov1998} develop a trust-region framework for managing approximation models. The central idea is to judge the approximation where it is being used and control the search accordingly. \citet{Sahigara2012} compare applicability-domain methods for quantitative structure--activity models. Their application differs from treatment modeling, but it supports the general need to distinguish nearby training support from distant prediction requests.

Distance restrictions and projection size answer different questions. An input can be near the training observations while its raw output requires a large projection. Conversely, a distant input can produce an apparently admissible state. Both indicators can be useful screening criteria, but neither is a calibrated probability of prediction accuracy. \citet{Hameed2026} develop trust-region filter methods for gray-box optimization with derivative information and feasibility considerations. Their contribution reinforces the need to separate approximation agreement, feasibility, and progress in the original objective.

Direct optimization provides a meaningful comparator only when its numerical formulation is explicit. \citet{BuzziFerraris2013} explain nonlinear systems and optimization in chemical engineering. \citet{Kraft1988} develop sequential quadratic programming, while \citet{WachterBiegler2006} develop an interior-point filter line-search method for nonlinear programming. These approaches support different numerical routes for the operating problem. A local termination condition does not establish the global optimum of a nonconvex plant model.

The underlying steady-state solution also needs attention. \citet{CurtissHirschfelder1952} establish the problem of stiff differential equations. Biological treatment contains processes with different response scales, which can make steady-state initialization difficult. \citet{KelleyKeyes1998} analyze pseudo-transient continuation as a route toward a steady solution. Time relaxation followed by an algebraic solve can provide a better initial state, but acceptance still depends on checking the original equations and component bounds.

\citet{AudetDennis2006} develop mesh adaptive direct search for constrained optimization. \citet{Ragonneau2022} study model-based derivative-free methods. These works explain why an operating search can remain possible when the projected prediction map is nonsmooth or reliable derivatives are unavailable. The comparison with a smooth direct formulation should account for the differences in numerical structure and stopping rules.

\citet{Bliek2023} benchmark surrogate optimization algorithms on expensive black-box functions. Their contribution makes assessment design and evaluation cost central to a performance comparison. \citet{Pedrozo2025} compare surrogate optimization with a reference process optimization setting. Both support evaluating the objective and constraints at the selected decisions using the original model. Search time, development time, and verification time should be reported separately because they answer different cost questions.

\citet{Zhang2026} develop an agile benchmarking framework for wastewater resource recovery technologies. Their emphasis on comparable assessment strengthens the need for common engineering priorities and clear process boundaries. Route-specific numerical scales can help solve the equations, but both routes must ultimately be judged using the same physical outputs and objective definition.

\section{The integrated research gap}

The literature provides strong foundations for component modeling, nonlinear statistical prediction, physical constraint enforcement, interpretable regression, and surrogate-assisted optimization. It also shows why these foundations cannot be reduced to one error score. A component surrogate must represent the response that is being assessed. A projection must enforce the relations that are being claimed. An interpretation must describe the prediction map that is actually returned. An operating decision must be evaluated against the original engineering problem.

The research gap concerns their integration in activated sludge analysis. A common reactor benchmark can compare raw and projected predictions across learning families. A structured regression can make influent, operating, and coupling terms inspectable while preserving a separate feasibility procedure. A connected plant formulation can extend enforcement to mixing, treatment stages, clarification, and solids inventory. Independent mechanistic evaluation can then assess the selected controls. This sequence supplies a coherent basis for studying predictive accuracy, physical admissibility, interpretability, and decision quality without treating any one of them as evidence for all the others.
